\documentclass{beamer}
\mode<presentation>
\usepackage{amsmath,amssymb,mathtools}
\usepackage{textcomp}
\usepackage{gensymb}
\usepackage{adjustbox}
\usepackage{subcaption}
\usepackage{enumitem}
\usepackage{multicol}
\usepackage{listings}
\usepackage{url}
\usepackage{graphicx} % <-- needed for images
\usepackage{xcolor}
\def\UrlBreaks{\do\/\do-}

\usetheme{Boadilla}
\usecolortheme{lily}
\setbeamertemplate{footline}{
	\leavevmode%
	\hbox{%
		\begin{beamercolorbox}[wd=\paperwidth,ht=2ex,dp=1ex,right]{author in head/foot}%
			\insertframenumber{} / \inserttotalframenumber\hspace*{2ex}
	\end{beamercolorbox}}%
	\vskip0pt%
}
\setbeamertemplate{navigation symbols}{}

\lstset{
	frame=single,
	breaklines=true,
	columns=fullflexible,
	basicstyle=\ttfamily\tiny    % tiny font so code fits
}

\numberwithin{equation}{section}

% ---- your macros ----
\providecommand{\nCr}[2]{\,^{#1}C_{#2}}
\providecommand{\nPr}[2]{\,^{#1}P_{#2}}
\providecommand{\mbf}{\mathbf}
\providecommand{\pr}[1]{\ensuremath{\Pr\left(#1\right)}}
\providecommand{\qfunc}[1]{\ensuremath{Q\left(#1\right)}}
\providecommand{\sbrak}[1]{\ensuremath{{}\left[#1\right]}}
\providecommand{\lsbrak}[1]{\ensuremath{{}\left[#1\right.}}
\providecommand{\rsbrak}[1]{\ensuremath{\left.#1\right]}}
\providecommand{\brak}[1]{\ensuremath{\left(#1\right)}}
\providecommand{\lbrak}[1]{\ensuremath{\left(#1\right.}}
\providecommand{\rbrak}[1]{\ensuremath{\left.#1\right)}}
\providecommand{\cbrak}[1]{\ensuremath{\left\{#1\right\}}}
\providecommand{\lcbrak}[1]{\ensuremath{\left\{#1\right.}}
\providecommand{\rcbrak}[1]{\ensuremath{\left.#1\right\}}}
\theoremstyle{remark}
\newtheorem{rem}{Remark}
\newcommand{\sgn}{\mathop{\mathrm{sgn}}}
\providecommand{\abs}[1]{\left\vert#1\right\vert}
\providecommand{\res}[1]{\Res\displaylimits_{#1}}
\providecommand{\norm}[1]{\lVert#1\rVert}
\providecommand{\mtx}[1]{\mathbf{#1}}
\providecommand{\mean}[1]{E\left[ #1 \right]}
\providecommand{\fourier}{\overset{\mathcal{F}}{ \rightleftharpoons}}
\providecommand{\system}{\overset{\mathcal{H}}{ \longleftrightarrow}}
\providecommand{\dec}[2]{\ensuremath{\overset{#1}{\underset{#2}{\gtrless}}}}
\newcommand{\myvec}[1]{\ensuremath{\begin{pmatrix}#1\end{pmatrix}}}
\newcommand{\mydet}[1]{\ensuremath{\begin{vmatrix}#1\end{vmatrix}}}

\newenvironment{amatrix}[1]{%
	\left[\begin{array}{@{}*{#1}{c}|c@{}}
	}{%
	\end{array}\right]
}

\newcommand{\myaugvec}[2]{\ensuremath{\begin{amatrix}{#1}#2\end{amatrix}}}
\let\vec\mathbf
% ---------------------

\title{12.263}
\author{EE25BTECH11042 - Nipun Dasari}

\begin{document}
	
	\begin{frame}
		\titlepage
	\end{frame}
	
	\begin{frame}{Question}
		The matrix
		\begin{align}
			\vec{A} = \begin{myvec}{1 & 1 & 1 \\ 2 & 1 & 2 \\ 1 & 3 & 2}\end{myvec}
		\end{align}
		can be decomposed into the product of a lower triangular matrix $\vec{L}$ and an upper triangular matrix $\vec{U}$ as $\vec{A} = \vec{L}\vec{U}$. Let $\vec{x}, \vec{z} \in \mathbb{R}^3$ and $\vec{b} = (1 \;\; 1 \;\; 1)^\top$.
		
		a) The solution $\vec{z} = (z_1 \;\; z_2 \;\; z_3)^\top$ of the system $\vec{L}\vec{z} = \vec{b}$ is
		\begin{enumerate}[label=\roman*)]
			\item $(-1 \;\; -1 \;\; -2)^\top$
			\item $(1 \;\; -1 \;\; 2)^\top$
			\item $(1 \;\; -1 \;\; -2)^\top$
			\item $(-1 \;\; 1 \;\; 2)^\top$
		\end{enumerate}
		
		b) The solution $\vec{x} = (x_1 \;\; x_2 \;\; x_3)^\top$ of the system $\vec{U}\vec{x} = \vec{z}$ is
		\begin{enumerate}[label=\roman*)]
			\item $(2 \;\; 1 \;\; -2)^\top$
			\item $(2 \;\; 1 \;\; 2)^\top$
			\item $(-2 \;\; -1 \;\; -2)^\top$
			\item $(-2 \;\; 1 \;\; -2)^\top$
		\end{enumerate}
	\end{frame}
	
	\begin{frame}{Theoretical solution}
		We will perform row operations on $\vec{A}$ to obtain $\vec{U}$. The multipliers used will form the entries of $\vec{L}$.
		\begin{align}
			\vec{A} = \begin{myvec}{1 & 1 & 1 \\ 2 & 1 & 2 \\ 1 & 3 & 2}\end{myvec}
		\end{align}
		Operation $R_2 \to R_2 - (\textcolor{red}{2})R_1$. The multiplier is $l_{21} = 2$.
		\begin{align}
			\begin{myvec}{1 & 1 & 1 \\ 0 & -1 & 0 \\ 1 & 3 & 2}\end{myvec}
		\end{align}
		Operation $R_3 \to R_3 - (\textcolor{red}{1})R_1$. The multiplier is $l_{31} = 1$.
		\begin{align}
			\begin{myvec}{1 & 1 & 1 \\ 0 & -1 & 0 \\ 0 & 2 & 1}\end{myvec}
		\end{align}
		\end{frame}
			
			\begin{frame}{Theoretical solution}
				Operation $R_3 \to R_3 - (\textcolor{red}{-2})R_2$. The multiplier is $l_{32} = -2$.
				\begin{align}
					\vec{U} = \begin{myvec}{1 & 1 & 1 \\ 0 & -1 & 0 \\ 0 & 0 & 1}\end{myvec}
				\end{align}
				The lower triangular matrix $\vec{L}$ is constructed from the multipliers.
				\begin{align}
					\vec{L} = \begin{myvec}{1 & 0 & 0 \\ l_{21} & 1 & 0 \\ l_{31} & l_{32} & 1}\end{myvec} = \begin{myvec}{1 & 0 & 0 \\ 2 & 1 & 0 \\ 1 & -2 & 1}\end{myvec}
				\end{align}
			\end{frame}
			
			\begin{frame}{Theoretical solution}
				Part a): Solving $\vec{L}\vec{z} = \vec{b}$
				
				We solve the system using Gaussian elimination.
				\begin{align}
					\myaugvec{3}{1 & 0 & 0 & 1 \\ 2 & 1 & 0 & 1 \\ 1 & -2 & 1 & 1} 
					&\xleftrightarrow[R_3 \to R_3 - R_1]{R_2 \to R_2 - 2R_1} \myaugvec{3}{1 & 0 & 0 & 1 \\ 0 & 1 & 0 & -1 \\ 0 & -2 & 1 & 0} \\
					&\xleftrightarrow{R_3 \to R_3 + 2R_2} \myaugvec{3}{1 & 0 & 0 & 1 \\ 0 & 1 & 0 & -1 \\ 0 & 0 & 1 & -2}
				\end{align}
				The solution is $\vec{z} = (1 \;\; -1 \;\; -2)^\top$, which is option (iii).
			\end{frame}
			
			\begin{frame}{Theoretical solution}
				Part b): Solving $\vec{U}\vec{x} = \vec{z}$
				
				We solve the system using the augmented matrix.
				\begin{align}
					\myaugvec{3}{1 & 1 & 1 & 1 \\ 0 & -1 & 0 & -1 \\ 0 & 0 & 1 & -2}
					&\xleftrightarrow{R_2 \to -R_2} \myaugvec{3}{1 & 1 & 1 & 1 \\ 0 & 1 & 0 & 1 \\ 0 & 0 & 1 & -2} \\
					&\xleftrightarrow[R_1 \to R_1 - R_2]{R_1 \to R_1 - R_3} \myaugvec{3}{1 & 0 & 0 & 2 \\ 0 & 1 & 0 & 1 \\ 0 & 0 & 1 & -2}
				\end{align}
				The solution is $\vec{x} = (2 \;\; 1 \;\; -2)^\top$, which is option (i).
			\end{frame}
			
			\begin{frame}{Plot}
				\begin{figure}[h!]
					\centering
					\includegraphics[width=0.6\columnwidth]{figs/Figure_1.png} 
					\caption*{Fig }
					\label{Fig1}
				\end{figure}
			\end{frame}
			
		\end{document}